
\documentclass[main]{subfiles} 

\begin{document}

\chapter{\MakeUppercase{REQUIREMENT ANALYSIS}}

Requirement analysis is a crucial phase of the project development process that identifies the technical, functional, and operational needs of the proposed system. This chapter outlines the system requirements, functional and non-functional requirements, and the feasibility of implementing the proposed temporal machine learning framework for predicting urban growth and LULC dynamics using Landsat time-series data.

\section{Requirement Analysis}
The proposed system is a data-driven analytical framework designed to process multi-temporal Landsat imagery, generate LULC maps, and apply temporal machine learning models to forecast future urban growth and LULC changes. The system requires adequate computational resources, appropriate software tools, and a stable execution environment to support data preprocessing, model training, evaluation, and visualization.

\subsection{Software Requirements}
The software requirements focus on open-source tools and platforms suitable for remote sensing analysis and machine learning.

\begin{enumerate}[label=\roman*.]
    \item \textbf{Operating System}
    \begin{itemize}
        \item Windows 11 / Linux / macOS
    \end{itemize}
    
    \item \textbf{Programming Language}
    \begin{itemize}
        \item Python (version 3.8 or higher)
    \end{itemize}
    
    \item \textbf{Remote Sensing and GIS Tools}
    \begin{itemize}
        \item Google Earth Engine (GEE) for accessing and preprocessing Landsat imagery
        \item QGIS or ArcGIS for spatial visualization and map validation
    \end{itemize}
    
    \item \textbf{Machine Learning and Deep Learning Libraries}
    \begin{itemize}
        \item TensorFlow or PyTorch for implementing temporal machine learning models
        \item Keras (if using TensorFlow backend) for rapid model prototyping
        \item Scikit-learn for baseline models and evaluation metrics
    \end{itemize}
    
    \item \textbf{Data Processing and Visualization Libraries}
    \begin{itemize}
        \item NumPy and Pandas for numerical and tabular data processing
        \item Matplotlib and Seaborns for plotting and result visualization
    \end{itemize}
    
    \item \textbf{Development Environment}
    \begin{itemize}
        \item Jupyter Notebook or Visual Studio Code for development and experimentation
    \end{itemize}
\end{enumerate}

\section{Functional Requirements}
Functional requirements describe what the system is expected to do. These requirements define the core functionalities necessary to achieve the project objectives.

\begin{enumerate}[label=\roman*.]
    \item The system shall acquire multi-year Landsat imagery for the selected study area.
    \item The system shall preprocess satellite images, including cloud masking, compositing, and normalization.
    \item The system shall generate annual LULC maps from Landsat data using suitable classification methods.
    \item The system shall construct temporal sequences of LULC data for machine learning input.
    \item The system shall implement temporal machine learning model.
    \item The system shall forecast future urban growth and LULC patterns over defined time horizons.
    \item The system shall evaluate model performance using accuracy metrics and hindcasting techniques.
    \item The system shall visualize historical and predicted LULC maps for interpretation and analysis.
    \item The system shall store trained models and prediction outputs for further analysis and reporting.
\end{enumerate}

\section{Non-Functional Requirements}
Non-functional requirements describe the quality attributes and operational constraints of the system.

\begin{enumerate}[label=\roman*.]
    \item \textbf{Performance}: The system should efficiently process large raster datasets, and model training should be feasible on a personal computer.
    
    \item \textbf{Scalability}: The framework should support extension to other geographic regions and additional years of Landsat data with minimal modification.
    
    \item \textbf{Reliability}: The system should produce consistent results under identical conditions, with preprocessing steps minimizing errors from cloud cover and noise.
    
    \item \textbf{Usability}: The workflow should be understandable and reproducible for users with basic GIS and Python knowledge, with well-documented code.
    
    \item \textbf{Maintainability}: The system should be modular, allowing independent updates to preprocessing, modeling, and evaluation components.
\end{enumerate}

% \newpage
\section{Feasibility Study}
The feasibility study evaluates whether the proposed project can be realistically implemented within the available constraints of time, resources, and technical capability.

\subsection{Technical Feasibility}
The project is technically feasible due to the availability of long-term Landsat data and mature machine learning frameworks. Google Earth Engine provides seamless access to preprocessed Landsat datasets, reducing data acquisition and preprocessing complexity. Open-source libraries such as TensorFlow and PyTorch enable the implementation of advanced temporal machine learning models without licensing constraints.

\subsection{Economic Feasibility}
The project is economically feasible as it relies entirely on freely available satellite data and open-source software tools. No additional cost is required for data acquisition or software licenses, making the approach suitable for academic and research-oriented applications.

\subsection{Operational Feasibility}
The system can be operated by users with basic knowledge of remote sensing, GIS, and Python programming. The workflow does not require specialized hardware beyond a standard personal computer, ensuring ease of deployment and operation in academic environments.

\subsection{Schedule Feasibility}
The project is feasible within the allocated time frame of a minor project. The modular structure of the workflow allows parallel progress in data preprocessing, model development, and evaluation, ensuring timely completion.

\end{document}
